\section{Глава 4. Динамическое программирование}
\subsection{Упражнение 4.1(c. 105)}
В примере 4.1, если $\pi$ – равновероятная случайная стратегия, чему равно $q_{\pi}(11, \text{down})$? А чему равно $q_{\pi}(7, \text{down})$?
\begin{solution}
    \[
    q_{\pi}(s,a)=\sum_{s',r} p(s',r\mid s,a)\,\bigl[r+\gamma v_{\pi}(s')\bigr]
    \]
    Из постановки задачи 
    \[\gamma=1,\quad r=-1\ \text{на всех переходах},\quad v_{\pi}(T)=0.\]
    \[
    q_{\pi}(11,\text{down})
    =1\cdot\,[-1+v_{\pi}(T)]=-1 + 0 = -1.
    \]
    
    \[
    q_{\pi}(7,\text{down})
    =1\cdot\,[-1+v_{\pi}(11)]
    =-1-14=-15.
    \]
\end{solution}

\subsection{Упражнение 4.2(c. 106)}
В примере $4.1$ предположим, что добавлено новое состояние $15$ под состоянием $13$, а действия left, up, right, down в нем переводят агента в состояния $12, 13, 14$ и $15$ соответственно. Переходы из старых состояний не изменились.
Чему тогда равно $v_{\pi}(15)$ при следовании равновероятной случайной стратегии? Теперь предположим, что динамика состояния $13$ тоже изменилась, так что действие down переводит агента в новое состояние $15$. Чему в этом случае равно $v_{\pi}(15)$ при следовании равновероятной случайной стратегии? 
\begin{solution}
    Уравнение Беллмана:
    \[
    v_{\pi}(s)=\sum_a \pi(a\mid s)
    \sum_{s'} p(s'\mid s,a)\,[ -1 + v_{\pi}(s') ]
    \]
    
    Рассмотрим \textbf{случай без изменения переходов}. Тогда
    \[
    v_{\pi}(15)
    =\frac14\Big[
    (-1+v_{\pi}(12))
    +(-1+v_{\pi}(13))
    +(-1+v_{\pi}(14))
    +(-1+v_{\pi}(15))
    \Big].
    \]
    
    Отсюда
    \[
    4v_{\pi}(15)
    =
    -4+v_{\pi}(12)+v_{\pi}(13)+v_{\pi}(14)+v_{\pi}(15),
    \]
    \[
    3v_{\pi}(15)
    =
    -4+v_{\pi}(12)+v_{\pi}(13)+v_{\pi}(14).
    \]
    
    Из примера 4.1:
    \[
    v_{\pi}(12)=-22,\qquad
    v_{\pi}(13)=-20,\qquad
    v_{\pi}(14)=-14.
    \]
    
    Итого,
    \[
    v_{\pi}(15)
    =\frac{-4-22-20-14}{3}
    =-20.
    \]

    Перейдем к случаю, в котором \textbf{изменения переходов произошли}. Из математической теории уравений Беллмана известно, что в случае, когда стратегия $\pi$ гарантирует завершение из любого состояния, решение системы относительно $v_{\pi}$ существует и единственно. Поэтому нам достаточно привести пример, при котором будет выполняться система - это и есть решение. Заметим: если мы присвоим $v_{\pi}(15) = v_{\pi}(13) = 20$, то  $v_{\pi}(13)$ останется равным $20$ и после изменения переходов(так как до изменений действие down из $13$ возвращало состояние с ценностью $20$). Более того, уравнение Беллмана будет выполняться для $v_{\pi}(15) = 20$ на основании первой части решения задачи. Значения ценностей остальных состояний не изменятся.
\end{solution}

\subsection{Упражнение 4.3(c. 107)}
Как выглядят уравнения, аналогичные $(4.3)$, $(4.4)$ и  $(4.5)$, для функции ценности действий $q_{\pi}$ и ее последовательных приближений функциями $q_0, q_1, q_2,\dots$
\begin{align}
v_{\pi}(s)
&= \mathbb{E}_{\pi}\!\left[ G_t \mid S_t = s \right] \notag\\
&= \mathbb{E}_{\pi}\!\left[ R_{t+1} + \gamma G_{t+1} \mid S_t = s \right] \notag \\
&= \mathbb{E}_{\pi}\!\left[ R_{t+1} + \gamma v_{\pi}(S_{t+1}) \mid S_t = s \right] \tag{4.3} \\
&= \sum_{a} \pi(a \mid s)
\sum_{s',r} p(s', r \mid s, a)
\left[ r + \gamma v_{\pi}(s') \right]. \tag{4.4}
\end{align}
\begin{align}
v_{k+1}(s)
&= \mathbb{E}_{\pi}\!\left[ R_{t+1} + \gamma v_k(S_{t+1}) \mid S_t = s \right] \notag\\
&= \sum_{a} \pi(a \mid s)
\sum_{s',r} p(s', r \mid s, a)
\left[ r + \gamma v_k(s') \right].
\tag{4.5}
\end{align}
\begin{solution}
    \begin{align}
    q_{\pi}(s,a)
    &= \mathbb{E}_{\pi}\!\left[ G_t \mid S_t = s, A_t = a \right] \notag\\
    &= \mathbb{E}_{\pi}\!\left[ R_{t+1} + \gamma G_{t+1}
    \mid S_t = s, A_t = a \right] \notag\\
    &= \mathbb{E}_{\pi}\!\left[
    R_{t+1} + \gamma
    \sum_{a'} \pi(a' \mid S_{t+1})
    q_{\pi}(S_{t+1},a')
    \mid S_t = s, A_t = a
    \right] \notag\\
    &=
    \sum_{s',r} p(s', r \mid s, a)
    \left[
    r + \gamma
    \sum_{a'} \pi(a' \mid s')
    q_{\pi}(s',a')
    \right].
    \notag
    \end{align}
    
    \begin{align}
    q_{k+1}(s,a)
    &=
    \mathbb{E}_{\pi}\!\left[
    R_{t+1}
    +
    \gamma
    \sum_{a'} \pi(a' \mid S_{t+1})
    q_k(S_{t+1},a')
    \mid S_t = s, A_t = a
    \right] \notag \\
    &=
    \sum_{s',r} p(s', r \mid s, a)
    \left[
    r + \gamma
    \sum_{a'\in A(s')} \pi(a' \mid s')
    q_k(s',a')
    \right] \notag
    \end{align}
\end{solution}

\subsection{Упражнение 4.4(c. 111)}
В алгоритме итерации по стратегиям на стр. $110$ есть тонкая ошибка, состоящая в том, что он может никогда не завершиться, если стратегия постоянно переключается между двумя или более стратегиями, которые одинаково хороши. В педагогических целях это, может, и неплохо, но на практике не годится. Модифицируйте псевдокод, так чтобы сходимость была гарантирована.
\begin{figure}[H]
    \centering
    \includegraphics[width=1\linewidth]{iter_alg.png}
    \caption{}
\end{figure}
\begin{solution}
    Проблема возникает на этапе улучшения стратегии. Если множество $A'(s) = \text{Arg} \max_{a \in A(s)}
    q_{\pi}(s,a)$  состоит более чем из одного элемента и содержит $\pi(s)$, то нерегламентированный выбор действия ведет к появлению циклов. \\
    Поэтому в процесс улучшения стратегии следует добавить следующую проверку:
    \[
    \text{Если} \ \pi(s) \in A'(s), \ \text{то} \ \pi(s) \leftarrow \pi(s) \]
    \[
    \text{В противном случае } \pi(s) \leftarrow \forall a \in A'(s)
    \]
\end{solution}

\subsection{Упражнение 4.5(c. 111)}
Как следовало бы определить алгоритм итерации по стратегиям
для ценностей действий? Приведите полный алгоритм вычисления $q_*$, аналогичный показанному на стр. 110 для вычисления $v_*$. Обратите особое внимание на это упражнение, потому что развиваемые в нем идеи будут использоваться в книге далее. 
\begin{solution}
    Оценивание стратегии в терминах ценности действий приведено в упражнении 4.3. \\
    Заметим, что обновление стратегии по правилу
    \[\pi'(s) = \arg\max_aq_{\pi}(s, a)\]
    при известной ценности действий еще проще и не требует дополнительных равенств. \\
    Надо показать, что имеет место строгое улучшение стратегий:
    рассмотрим две детерминированные стратегии $\pi$ и $\pi'$ такие, что \[\pi'(s) = \arg\max_aq_{\pi}(s, a)\] 
    Предположим, что улучшения политики не произошло, тогда $q_{\pi'}(s,a) = q_{\pi}(s,a)$.
    Получим
    \[\pi'(s) = \arg\max_aq_{\pi}(s, a) = 
    \arg\max_aq_{\pi'}(s, a) \]
    Запишем уравнение Беллмана:
    \begin{align*}
        q_{\pi'}(s,a) = \\
        =& \sum_{s'\in S}\sum_{r\in R} p(s',r\mid s,a)\,\Bigl(r + \gamma\sum_{a'\in A(s')} \pi'(a'\mid s')\, q_{\pi'}(s',a')\Bigr) =\\
        =& \sum_{s'\in S}\sum_{r\in R} p(s',r\mid s,a)\,\Bigl(r + \gamma \cdot q_{\pi'}(s',\pi'(s'))\Bigr) =\\
        =& \sum_{s'\in S}\sum_{r\in R} p(s',r\mid s,a)\,\Bigl(r + \gamma \cdot \max_{a'} q_{\pi'}(s',a')\Bigr)
    \end{align*}
    Крайние равенства соответствуют уравнению оптимальности Беллмана, поэтому $q_{\pi'}(s,a) = q_{\pi}(s,a) = q_*(s,a)$ и $\pi(s) = \pi'(s) = \pi_*(s)$.
    \vspace{5mm} \\
    Псевдокод алгоритма: \\
    \textbf{Вход:} множества $S,\ A(s)$, динамика $p(s',r\mid s,a)$, $\gamma\in[0,1)$, порог $\theta>0$.\\
    \textbf{Выход:} $q_*$ и оптимальная стратегия $\pi_*$.
    
    \textbf{1. Инициализация}\\
    Для всех $s\in S$ выбрать $\pi(s)\in A(s)$ произвольно.\\
    Задать $q(s,a)$ произвольно для всех $(s,a)$.
    
    \textbf{2. Оценивание стратегии (итеративное)}\\
    \textbf{Повторять:}\\
    \hspace*{5mm} $\Delta \leftarrow 0$\\
    \hspace*{5mm} \textbf{для каждого } $s\in S$ \textbf{и каждого } $a\in A(s)$:\\
    \hspace*{5mm}\hspace*{5mm} $old \leftarrow q(s,a)$\\
    \hspace*{5mm}\hspace*{5mm} $q(s,a)\leftarrow 
    \sum_{s',r} p(s',r\mid s,a)\Bigl[r+\gamma\, q\bigl(s',\pi(s')\bigr)\Bigr]$\\
    \hspace*{5mm}\hspace*{5mm} $\Delta \leftarrow \max\bigl(\Delta,\ |old-q(s,a)|\bigr)$\\
    \textbf{пока } $\Delta \ge \theta$
    
    \textbf{3. Улучшение стратегии}\\
    $policy\_stable \leftarrow true$\\
    \textbf{для каждого } $s\in S$:\\
    \hspace*{5mm} $old\_action \leftarrow \pi(s)$\\
    \hspace*{5mm} $\pi(s)\leftarrow \arg\max_{a\in A(s)} q(s,a)$\\
    \hspace*{5mm} \textbf{если } $old\_action \ne \pi(s)$, \textbf{то } $policy\_stable \leftarrow false$\\
    \textbf{если } $policy\_stable$, \textbf{то остановиться и вернуть} $q, \ \pi$\\
    \textbf{иначе перейти к 2.}
\end{solution}

\subsection{Упражнение 4.6(c. 112)}
Предположим, что мы ограничены только $\varepsilon$-мягкими стратегиями, для которых вероятность выбора любого действия в каждом состоянии $s$ не меньше $\frac{\varepsilon}{|A(s)|}$. Качественно опишите изменения, которые нужно было бы внести
на шагах 3, 2 и 1 (именно в таком порядке) приведенного на стр. 110 алгоритма итерации по стратегиям для вычисления $v_*$.
\begin{solution}

    \textbf{3. Улучшение стратегии}\\
    policy-stable $\leftarrow$ true\\
    Для каждого $s \in S$:\\
    \hspace*{5mm} old-action $\leftarrow \arg\max_{a\in A(s)} \pi(a\mid s)$\\
    \hspace*{5mm} $a_* \leftarrow \arg\max_{a\in A(s)}
    \sum_{s',r} p(s',r\mid s,a)\bigl[r+\gamma V(s')\bigr]$\\
    \hspace*{5mm} Для каждого $a \in A(s)$:\\
    \hspace*{5mm}\hspace*{5mm} если $a=a_*$
    то 
    $\pi(a\mid s) \leftarrow 1-\varepsilon+\dfrac{\varepsilon}{|A(s)|}$\\
    \hspace*{5mm}\hspace*{5mm} иначе 
    $\pi(a\mid s) \leftarrow \dfrac{\varepsilon}{|A(s)|}$\\
    \hspace*{5mm} если old-action $\neq a_*$ то policy-stable $\leftarrow$ false\\
    Если policy-stable, то остановиться и вернуть $v_*$,$\pi_*$;
    иначе перейти к 2.
    
    \vspace{5mm}
    
    \textbf{2. Оценивание стратегии}\\
    Повторять: \\
    \hspace*{5mm} $\Delta \leftarrow 0$\\
    \hspace*{5mm} Повторять для каждого $s \in S$: \\
    \hspace*{5mm}\hspace*{5mm} $v \leftarrow V(s)$\\
    \hspace*{5mm}\hspace*{5mm} $V(s) \leftarrow 
    \sum_{a \in A(s)} \pi(a \mid s)
    \sum_{s',r} p(s', r \mid s, a)\,[r + \gamma V(s')]$\\
    \hspace*{5mm}\hspace*{5mm} $\Delta \leftarrow \max(\Delta,\ |v - V(s)|)$\\
    пока не окажется $\Delta < \theta$

    \vspace{5mm}

    \textbf{1. Инициализация}\\
    $V(s) \in \mathbb{R} \text{ выбираются произвольно для всех } s \in S$\\
    $\pi(a \mid s) \text{ выбирается произвольно при условии }
    \pi(a \mid s) \ge \dfrac{\varepsilon}{|A(s)|}
    \text{ для всех } s \in S \text{ и } a \in A(s)$
\end{solution}

\subsection{Упражнение 4.7(c. 112)}
Напишите программу, реализующую итерацию по стратегиям, и решите задачу об аренде машин со следующими изменениями. Один из служащих Джека в первом офисе каждую ночь ездит на автобусе домой, а живет он рядом со вторым офисом. Он с удовольствием готов
перегонять одну машину во второй офис бесплатно. Но каждая дополнительная машина все равно стоит 2 доллара, равно как и перегон всех машин в обратном направлении. Кроме того, в  каждом офисе у  Джека ограниченное количество парковочных мест. Если на ночь остается больше $10$ машин (после всех перегонов), то за пользование другой парковкой взимается дополнительно $4$ доллара
(вне зависимости от того, сколько машин там стоит). Такого рода нелинейности и  произвольная динамика часто встречаются в  реальных задачах, и  все методы оптимизации, кроме динамического программирования, к ним применимы с трудом. Для проверки своей программы сначала воспроизведите результаты
для исходной задачи.
\begin{solution}
    \textbf{Код смотри в отдельном файле репозитория.} \\
    Основную трудность составляет вычисление модели окружающей среды $p(s', r \mid s, a)$. Для более эффективных вычислений разобьем на два шага:
    \begin{equation}
    G(b_1,b_2) = \sum_{r_1,r_2}\mathbb{P}(r_1)\mathbb{P}(r_2)V(b_1+r_1,b_2+r_2)
    \tag{1}
    \end{equation}
    В этой формуле $b_1, b_2$ - ежедневное состояние офисов после запроса, но перед возвратом. Тогда $b_1 + r_1, b_2 + r_2$ - состояние офисов после возврата, где $r_1, r_2$ - всевозможные значения случайной величины возврата. Эта формула представляет собой математическо ожидание по возвратам(усредненная ценность) \\
    Рассмотрим второй шаг:
    \begin{equation}
    V(s_1,s_2) = \sum_{q_1,q_2}\mathbb{P}(q_1)\mathbb{P}(q_2)\left[r + \gamma G(b_1,b_2)\right]
    \tag{2}
    \end{equation}
    В этой формуле $s_1, s_2$ - состояние среды перед запросом, но после перегона, $q_1, q_2$ - запросы в первом и втором офисе соответственно, $b_1 = s_1 - q_1, b_2 = s_2 - q_2$ - ежедневное состояние офисов после запроса, но перед возвратом.
    По аналогии эта формула является усреднением ценности по запросам. Именно ей я пользуюсь в коде как при оценивании, так и при улучшении стратегии. \\
    В приложенном ноутбуке решается исходный пример, а после он модифицируется под упражнение. Некоторые области полученного графика не поддаются логическому объяснению. Я связываю это с произвольной и нелинейной динамикой.
    \begin{figure}[H]
    \centering

    \begin{subfigure}{0.48\textwidth}
        \centering
        \includegraphics[width=\linewidth]{fig_4_7_final.png}
        \caption{Решение упр.4.7}
    \end{subfigure}
    \hfill
    \begin{subfigure}{0.48\textwidth}
        \centering
        \includegraphics[width=\linewidth]{fig_4_7_sol.png}
        \caption{Решение исходного примера}
    \end{subfigure}

    \caption{Сравнение результатов}
    \end{figure}
\end{solution}

\subsection{Упражнение 4.8(c. 114)}
Почему у оптимальной политики в задаче игрока такая странная
форма? В частности, для капитала $50$ предлагается поставить все на одно подбрасывание, а для капитала $51$ - нет. Почему такая стратегия хороша? 
\begin{solution}
    Запишем формулу для оптимальной политики:
    \[\pi_*(s) = \arg\max_a\ \left[p_h\cdot V_*(s+a) + 
    (1 - p_h) \cdot V_*(s-a)\right]\]
    В условии задачи $p_h < 0.5$, поэтому игроку разумно действовать осторожно. При $s=51$ максимальная ставка равна $49$, но в случае неудачи, с вероятностью $0.6$, он окажется в состоянии $s=2$, ценность которого мала. Поэтому игрок выбирает поставить $1$ доллар и в худшем случае оказаться в состоянии $s=50$, которое является более оптимистичным.
\end{solution}

\subsection{Упражнение 4.9(c. 115)}
Реализуйте алгоритм итерации по ценности для задачи игрока и решите ее для $p_h = 0.25$ и $p_h= 0.55$. Возможно, вы сочтете удобным ввести два фиктивных состояния, соответствующих завершению с капиталом $0$ и $100$, сопоставив им ценности $0$ и $1$. Представьте свои результаты графически, как на рис.4.3. Является ли полученное решение устойчивым при 
$\theta \rightarrow  0$?
\begin{solution}
    \textbf{Код смотри в отдельном файле репозитория.} При $\theta \rightarrow  0$ оценивание политики сходится к единственной оптимальной функции ценности, но выбор действий в оптимальной политике определен не единственным образом(для сравнения результатов в коде я разрешал неоднозначности, выбирая минимальную ставку)
    \begin{figure}[H]
    \centering

    \begin{subfigure}{0.8\textwidth}
        \centering
        \includegraphics[width=\linewidth]{fig_4_9_sol_v.png}
        \caption{Оптимальная ценность состояний}
    \end{subfigure}
    \hfill
    \begin{subfigure}{0.8\textwidth}
        \centering
        \includegraphics[width=\linewidth]{fig_4_9_pi.png}
        \caption{Оптимальная политика}
    \end{subfigure}

    \caption{}
    \end{figure}
\end{solution}

\subsection{Упражнение 4.10(c. 115)}
Каков аналог обновления в  алгоритме итерации по ценности
$(4.10)$ для ценностей действий $q_{k+1}(s, a)$? 
\begin{align*}
    v_{k+1}(s) &= \\
    &=\max_a\mathbb{E}\left[ R_{t+1}+\gamma\cdot v_{k}(S_{t+1} )\mid S_t=s, A_t=a) \right] =\\
    &= \max_a\sum_{s',r}p(s',r\mid s,a)\cdot\left[r + \gamma\cdot v_k(s')\right]
\tag{4.10}
\end{align*}
\begin{solution}
    Как и в случае ценности состояний, эта формула следует либо из упражнения 4.5, либо из уравнения оптимальности Беллмана для ценности действий:
    \[
        q_{k+1}(s,a)
        =
        \sum_{s',r}
        p(s',r \mid s,a)
        \left[
        r + \gamma \max_{a'} q_k(s', a')
        \right]
    \]
\end{solution}
